\chapter{实验评估}
\label{chapter:evaluation}

本章通过一系列综合实验对所提出的隐私保护程序验证框架进行系统性评估。实验的核心目的是验证本文提出的两阶段协议——即从 SSA 中间表示到 SAT 公式的转换（阶段一），以及基于 Halo2 证明系统的 SAT/UNSAT 零知识验证（阶段二）——在实际应用场景中的性能表现。我们重点关注三个关键指标：执行效率、证明规模以及面对不同规模问题时的可扩展性。

具体的评估目标包括：
\begin{enumerate}
    \item \textbf{阶段一转换效率}：深入对比引入 HyperPlonk 优化策略的 zkWASM 与原生通用 zkVM 在程序到公式转换阶段的性能差异，分析优化方案在处理复杂逻辑时的加速效果。
    \item \textbf{阶段二验证性能}：
    \begin{itemize}
        \item 针对 UNSAT 实例，评估专用 Halo2 电路相对于通用 zkVM 的效率提升，特别是验证核心技术“挑战值编码”在处理变长子句时的有效性。
        \item 针对 SAT 实例，验证“双矩阵赋值”方案是否实现了与公式规模无关的常数级验证复杂度，以及其在并行约束检查中的表现。
    \end{itemize}
    \item \textbf{证明紧凑性}：量化评估协议生成的零知识证明在存储占用和网络传输上的开销优势，验证其在带宽受限环境下的适用性。
\end{enumerate}

\section{实验设置}

\subsection{硬件与软件环境}
为了确保实验数据的可靠性与一致性，所有实验均在统一的高性能计算工作站上进行。硬件配置方面，选用 AMD EPYC 7402 处理器（主频 2.80GHz，配备 32GB 内存），该配置能够模拟真实的服务器计算环境。软件环境基于 Ubuntu 22.04 LTS 操作系统。

在软件栈实现上，阶段一基于 zkWASM 开源框架进行定制修改，集成了 HyperPlonk 证明系统以优化底层的约束生成逻辑，特别是针对位运算和布尔逻辑的电路表达进行了改进。在阶段二中，密码学后端采用 Rust 语言实现的 Halo2 库，多项式承诺方案选用高效的 KZG（Kate-Zaverucha-Goldberg）方案，椭圆曲线参数选用以太坊生态通用的 BN254 曲线，以保证系统具备良好的互操作性。

\subsection{基准测试集构造方法}

为全面评估基于Halo2算术化的零知识可满足性证明方案，我们设计并实现了一套系统性的基准测试集构造流程。本实验的核心目标在于验证方案的两个关键特性：其验证开销应与可满足性赋值或归结证明的\textbf{宽度}（即子句中文字的最大数量）基本无关，同时与不可满足实例的归结证明\textbf{长度}（即推理步数）亦保持较低的关联性。为此，实验分为两个主要阶段。

在第一阶段，我们以业界广泛使用的Boolector基准测试套件作为原始输入源。该套件包含了从硬件验证问题转换而来的丰富SAT实例。对于每一个测试用例，我们首先利用基于zkWASM虚拟机实现的符号执行引擎与程序转换系统，对其进行处理。该系统能够解析程序的中间表示，通过路径探索与约束收集，动态地生成一个等价的布尔可满足性问题，并将其输出为标准合取范式（CNF）格式的SAT公式。此步骤确保了我们的验证流程始于真实的程序语义。

随后，我们使用轻量级但健壮的MiniSAT求解器对生成的CNF公式进行求解。根据公式的可满足性，求解器输出两类结果：对于可满足（SAT）的实例，输出一组使得整个公式为真的变量赋值；对于不可满足（UNSAT）的实例，则输出一个形式化的归结反驳证明，该证明记录了从初始子句导出空子句（矛盾）的每一步归结推理步骤。这两类输出构成了下一阶段验证的“见证”。我们将SAT实例的变量赋值系统地编码为一个二维的\textbf{可满足性赋值矩阵}；将UNSAT实例的归结证明序列编码为一个\textbf{归结证明追踪矩阵}。这两个矩阵的结构均被设计为可直接映射到Halo2零知识证明系统的电路约束中，从而实现对求解结果的简洁且可验证的表示。

为精确量化验证性能与关键参数的关系，我们通过受控的方法构造了三组具有特定特征的测试集：

\textbf{1. 宽度无关性测试集构造}
为隔离并测量子句宽度对验证性能的影响，我们采用“语义保持宽度扩展”技术。具体而言，对于一个给定的原始CNF公式，我们首先生成一个足够大的、与公式中所有现有变量都不相交的“新鲜变量”池。然后，对于公式中的每一个子句，我们根据目标宽度值，从其尾部开始，顺序地从新鲜变量池中取出正文字或负文字进行追加填充。通过精心设计填充变量的顺序和极性，可以确保新添加的文字不会与子句内或其他子句中的文字产生新的单元子句或导致子句有效性改变，从而严格保持公式的原始可满足性。通过系统性地调整填充量，我们从一个基础公式衍生出了一系列语义等价但最大子句宽度从5、10、20、50、100直至400的合成公式族。该构造方法使得宽度成为唯一主要变化的自变量，从而能够清晰观测验证时间随宽度增长的趋势。

\textbf{2. 长度无关性测试集构造}
为评估验证开销与逻辑推理链条长度的关系，我们专注于UNSAT实例。我们从Boolector基准中筛选出一批具有不同复杂度的UNSAT问题，并使用MiniSAT进行求解，同时要求其输出详细的归结证明。我们收集这些证明，并按其包含的归结步数（即证明长度）进行归类。最终得到的测试集覆盖了从简短证明（20-50步）、中等长度证明（200-500步）到超长证明（1000-2450步）的广泛范围。为进一步控制变量，对于某些能产生极长证明的“种子”公式，我们还通过调整求解器的启发式策略（如文字决策顺序）或对公式进行轻微的语义保持变形（如变量重命名），诱导求解器生成不同长度但证明同一结论的变体证明，以更纯粹地分析长度因素的影响。

\textbf{3. 综合规模与多样性测试集}
为检验系统在更一般场景下的鲁棒性，我们选取了一组在规模上具有高度多样性的原始基准。这些基准对应的SAT公式子句数量跨度从100条至11,248条，其转换前程序的静态单赋值形式节点数也从23个到2,537个不等。此集合包含了从简单的位向量运算到复杂的控制流与数据流交互的场景，旨在评估验证方案从单元级性质检验到中等规模集成程序验证的综合适应能力。

所有实验均在统一环境中运行，每个测试用例独立执行多次取平均时间，以消除随机误差。通过上述严谨的构造流程，我们能够分别分析验证时间与证明宽度、证明长度以及问题原始规模之间的函数关系，从而为所提出方案的效率与可扩展性提供确凿的数据支撑。

\section{基于zkWASM程序到公式的转换性能评估}

在评估的第一阶段，我们重点考察了整个隐私保护验证流程的初始环节——将静态单赋值形式的程序代码编译为等价布尔可满足性公式的转换效率。此阶段作为后续零知识验证的预处理步骤，其性能直接决定了系统处理实际问题的端到端延迟与资源消耗上限，是评估框架整体实用性的关键指标之一。为实现高效的转换，我们对标准的zkWASM（基于WebAssembly指令集的零知识虚拟机）证明协议进行了架构层面的优化，将其替换为基于HyperPlonk算术化方案构建的专用转换电路。HyperPlonk方案的核心优势在于其对高阶自定义逻辑门的原生支持，这允许我们将程序SSA表示中的复杂布尔操作（如异或、条件选择等）直接映射为紧凑的多项式约束集，从而在算术化层面实现更高效、更低开销的编码，避免了通用zkVM中因指令级模拟带来的额外间接开销。

表\ref{tab:phase1-perf}系统性地展示了九种不同结构复杂度与规模特征的程序样例在两种实现方案下的详细转换性能数据。实验结果表明，采用HyperPlonk优化的方案在所有测试用例中均实现了显著的性能提升，相对于原生zkWASM方案取得了平均约$3\times$的加速比。从具体数据来看，对于仅包含基本变量赋值操作的简单布尔程序，转换耗时从1.23秒降低至0.41秒，体现了基础操作开销的有效削减；而对于结构高度复杂、包含2,537个SSA节点的大规模综合性程序，原生方案需要长达27.40秒的转换时间，而优化方案仅需8.77秒，将转换延迟大幅缩短至原有水平的约三分之一，显著提升了系统面对复杂验证任务时的响应能力与用户交互体验。这种性能提升主要源于底层证明系统架构的深刻变革：原生zkWASM方案本质上是一个通用的虚拟机模拟器，它需要将SSA程序中的每一个逻辑运算指令首先翻译为等价的WebAssembly字节码，然后在该虚拟机的零知识执行环境中逐条模拟执行，并在此过程中为每一条指令的输入、输出及状态变迁生成轨迹约束，这一过程不可避免地引入了巨大的字节码编译开销、虚拟机解释循环开销以及为维持通用性而必需的大量冗余约束生成开销。相比之下，优化方案跳过了中间的虚拟机模拟层，通过利用HyperPlonk提供的灵活接口，将SSA数据流图的拓扑结构（例如，将多路选择器ITE结构映射为多路复用器门、将异或网络映射为自定义的XOR门）直接编译为一组预定义的高阶多项式约束。这种“直接编译”模式消除了通用执行轨迹中的大量中间变量和状态编码，显著降低了底层证明电路的逻辑门数量与互联复杂度，从而在根本上减少了多项式约束系统的规模与求解时间。进一步观察表中的数据趋势可以发现，无论是原生方案还是优化方案，其转换耗时均与程序规模的关键指标（如变量数量和生成的CNF子句数）呈现出近似线性的增长关系，这有力地证明了优化后的方案在显著提升绝对性能的同时，依然保持了优良的计算复杂度特性，在处理更大规模程序时不会出现性能的急剧劣化，具备良好的工程可扩展性。

\begin{table}[htbp]
\centering
\bicaption{ZKWASM与HyperPlonk增强版的转换性能对比}{Performance Comparison between ZKWASM and ZKWASM-hyperplonk}
\label{tab:phase1-perf}
\begin{tabular}{lccccc}
\toprule
\multirow{2}{*}{\textbf{程序类型}} & \multirow{2}{*}{\textbf{SSA节点}} & \multirow{2}{*}{\textbf{变量数}} & \multirow{2}{*}{\textbf{子句数}} & \multicolumn{2}{c}{\textbf{转换耗时 (s)}} \\
\cline{5-6}
& & & & \textbf{ZKWASM} & \textbf{Ours} \\
\midrule
布尔赋值 & 23 & 26 & 91 & 1.23 & 0.41 \\
与/或运算 & 34 & 47 & 203 & 2.45 & 0.78 \\
异或/非运算 & 27 & 41 & 175 & 2.23 & 0.71 \\
条件表达式 (ITE) & 38 & 73 & 285 & 2.87 & 0.92 \\
嵌套 ITE (深度3) & 64 & 127 & 498 & 4.87 & 1.56 \\
蕴含链 & 43 & 103 & 403 & 4.25 & 1.36 \\
混合布尔逻辑 & 89 & 231 & 960 & 8.05 & 2.58 \\
复杂断言 & 126 & 256 & 1,065 & 7.93 & 2.54 \\
大规模程序 & 2,537 & 2,719 & 11,248 & 27.40 & 8.77 \\
\bottomrule
\end{tabular}
\end{table}

\section{基于zkWASM的符号执行系统性能评估}

\section{基于Halo2的可满足性赋值与归结证明性能评估}

在评估的第二阶段，我们将核心关注点转向对已生成的布尔可满足性公式进行零知识验证的效率与特性分析。这是整个隐私保护验证框架的核心步骤，其性能直接决定了验证过程的实际可用性。根据公式的最终属性——即可满足或不可满足——系统需要执行两种在逻辑与算法上截然不同的验证流程。这两种流程均基于我们的核心优化方案构建，我们分别对它们的性能特征进行了详尽的量化评估。

\subsection{UNSAT实例的零知识归结证明验证}

对于被判定为不可满足的公式实例，零知识验证的核心任务在于确认由外部高效求解器（如SAT求解器）所提供的形式化归结证明的正确性，同时不泄露证明本身的具体细节。表\ref{tab:unsat-comparison} 系统性地对比了本文提出的基于Halo2前端与专用查找表电路的验证方案，与采用通用zkWASM虚拟机模拟验证逻辑的方案之间的性能差异。实验数据揭示了两者之间存在巨大的效率鸿沟。

通用zkVM方案因其设计目标是提供普适性的可验证计算环境，故在验证归结证明时，必须将每一步归结规则（如消解律的应用程序）编译并转化为一系列底层的WebAssembly虚拟机指令，并模拟其完整的执行轨迹以生成相应的约束。这种“模拟-证明”的模式导致了验证电路的复杂性与归结证明的步数呈强相关性，电路规模随步数增加而指数级膨胀。例如，当归结证明长度达到2,450步时，zkWASM方案需要构建规模高达$2^{22}$个约束的庞大电路，验证过程耗时长达81.90秒。

与此形成鲜明对比的是，我们基于Halo2构建的专用电路方案采用了完全不同的设计哲学。该方案的核心创新在于利用了Halo2算术化框架对查找表机制的高效原生支持。我们将归结证明验证的语义规则（例如，判定子句$C$是否合法地由子句$A$和$B$通过消解文字$l$推导得出）预先编码到一系列定制的查找表中。在验证时，每一步证明的验证操作被转化为一次或数次对相应查找表的高效查询与多项式承诺检查，从而完全避免了模拟整个验证程序状态机所带来的巨大开销。这种设计使得验证时间主要线性依赖于证明步数本身，而与底层模拟的复杂度脱钩。对于同样的2,450步复杂证明，本方案仅需2.80秒即可完成验证，相较于通用方案实现了近30倍的性能加速。这证明了为特定验证逻辑设计专用算术化电路在性能上的显著优势。

\begin{table}[htbp]
\centering
\bicaption{UNSAT验证性能对比：Halo2 vs ZKWASM}{ZK Resolution Proof Verification Time: Halo2 vs ZKWASM}
\label{tab:unsat-comparison}
\begin{tabular}{cccccc}
\toprule
\multirow{2}{*}{\textbf{归结长度}} & \multirow{2}{*}{\textbf{子句宽度}} & \multicolumn{2}{c}{\textbf{Halo2 (Ours)}} & \multicolumn{2}{c}{\textbf{ZKWASM}} \\
\cline{3-6}
 & & \textbf{电路规模} & \textbf{耗时} & \textbf{电路规模} & \textbf{耗时} \\
\midrule
20 & 5 & $2^6$ & 34.47 ms & $2^{18}$ & 1.20 s \\
274 & 34 & $2^9$ & 412.46 ms & $2^{19}$ & 8.33 s \\
1034 & 84 & $2^{11}$ & 1.50 s & $2^{21}$ & 42.54 s \\
1511 & 98 & $2^{11}$ & 1.60 s & $2^{21}$ & 43.78 s \\
2450 & 142 & $2^{12}$ & 2.80 s & $2^{22}$ & 81.90 s \\
\bottomrule
\end{tabular}
\end{table}

本工作的一个核心创新是提出了“挑战值编码”技术，旨在从根本上解决传统零知识验证方案在处理包含大量文字的宽子句时所面临的性能瓶颈。在经典的归结证明验证算法中，验证每一步归结操作的正确性通常需要对参与消解的子句进行逐文字的比较与扫描，其计算开销与子句的宽度（即子句中文字的数量）呈线性正相关，这在面对由程序分析生成的、往往包含宽子句的CNF公式时，会成为显著的效率制约因素。

为了严格且量化地验证我们所提出技术的有效性，我们设计了一组控制变量实验。具体而言，我们选取了多个具有不同归结证明长度的基础公式作为“种子”。对于每一个种子公式，我们通过系统性的语义保持变换——即在每个子句的尾部，顺序添加一组互不冲突且不影响公式可满足性的“填充”文字——来人工构造出一系列新的公式变体。这些变体与原始公式在逻辑上等价，具有完全相同的归结证明长度与结构，唯一的区别在于其子句的最大宽度被精确地控制为不同的目标值。通过这种方法，我们能够隔离“子句宽度”这一单一变量，并精确测量其对验证时间的影响。

表\ref{tab:width-independence} 的实验结果清晰地展示了验证开销对于子句宽度变化的强鲁棒性。数据显示，即使在子句宽度从80个文字大幅增加至400个文字（增长幅度达5倍）的情况下，对应证明的零知识验证时间几乎维持恒定，波动范围远低于测量误差允许的阈值。以证明长度为1034步的基准组为例，当宽度分别为80、150、200和400时，其验证耗时分别为1.49秒、1.53秒、1.48秒和1.51秒，始终稳定在1.5秒左右。在证明长度更长的2450步基准组中也观察到了完全一致的趋势。

这一实验结果提供了强有力的经验证据，证实了我们所提出的算法成功地实现了其设计目标：通过引入随机挑战值并对子句中所有文字的赋值进行线性组合，我们将原本需要对$W$个文字进行的、开销为$O(W)$的显式逐一扫描与比较操作，转化为对一个在有限域上计算的单一组合值的代数相等性检查，其开销为$O(1)$。这种转化在保持证明系统完备性与可靠性的前提下，从计算复杂度层面彻底解耦了验证开销与子句宽度之间的关联。

\begin{table}[htbp]
\centering
\bicaption{UNSAT验证的子句宽度无关性测试}{Clause Width-Independent Verification for UNSAT Cases}
\label{tab:width-independence}
\begin{tabular}{cccccc}
\toprule
\multicolumn{3}{c}{\textbf{基准组 1 (长度 1034)}} & \multicolumn{3}{c}{\textbf{基准组 2 (长度 2450)}} \\
\cmidrule(lr){1-3} \cmidrule(lr){4-6}
\textbf{公式长度} & \textbf{宽度} & \textbf{耗时} & \textbf{公式长度} & \textbf{宽度} & \textbf{耗时} \\
\midrule
1034 & 80 & 1.49 s & 2450 & 80 & 2.83 s \\
1034 & 150 & 1.53 s & 2450 & 150 & 2.85 s \\
1034 & 200 & 1.48 s & 2450 & 200 & 2.80 s \\
1034 & 400 & 1.51 s & 2450 & 400 & 2.91 s \\
\bottomrule
\end{tabular}
\end{table}

\subsection{SAT 实例的零知识可满足性赋值验证}

针对可满足的公式实例，零知识验证的目标是证明验证者掌握一组能够使整个合取范式公式为真的布尔变量赋值，同时确保该赋值内容本身不会泄露给证明验证方。为达成这一目标，我们创新性地设计了“双矩阵赋值”编码结构与相应的并行化验证电路。该结构将变量赋值信息与子句满足性检查逻辑分离，并利用多项式承诺技术实现高效的批量验证。

如表\ref{tab:sat-verify} 所示，该方案展现出卓越的验证效率稳定性。其证明电路规模被预先设定并恒定为 $2^8$ 行，这一规模主要由电路设计中用于并行处理的固定矩阵行数所决定，而与具体输入SAT公式的规模（包括变量总数和子句总数）完全无关。因此，无论待验证的SAT公式是仅包含498个子句的小型问题，还是扩展至包含4372个子句的中大型问题，系统完成整个零知识验证过程的总耗时均能稳定在400毫秒左右（实验测得范围在378.64 ms至417.88 ms之间）。

产生这种常数级别验证时间特性的根本原因在于方案内在的高度并行化约束检查机制。传统验证需要逐个检查每个子句是否被赋值满足，导致开销与子句数线性相关。而本方案通过将全部变量的赋值编码为一个多项式向量，并将所有子句的满足性条件统一表述为另一组多项式约束，从而能够在一个固定次数的多项式求值-验证过程中，一次性完成对所有子句的批量并行验证。这种设计巧妙地将原本与公式规模成线性甚至更高阶关系的计算复杂度，压缩为仅与电路固定参数相关的常数次域运算。这一特性使得本方案特别适用于需要频繁或实时验证大规模系统配置、硬件电路设计或复杂约束条件可满足性的应用场景。验证成本不会因问题复杂度的提升而发生显著增长，赋予了系统极高的可扩展性与实用性。

\begin{table}[htbp]
\centering
\bicaption{SAT赋值验证的常数级时间特性}{ZK SAT Assignment Verification Time}
\label{tab:sat-verify}
\begin{tabular}{cccc}
\toprule
\textbf{公式规模} & \textbf{子句宽度} & \textbf{验证耗时} & \textbf{电路规模} \\
\midrule
498 & 50 & 378.64 ms & $2^{8}$ \\
1065 & 50 & 395.27 ms & $2^{8}$ \\
4372 & 50 & 417.88 ms & $2^{8}$ \\
\bottomrule
\end{tabular}
\end{table}

\section{证明规模分析}

零知识证明的简洁性直接影响系统的存储成本。表\ref{tab:proof-size} 对比了 UNSAT 场景下的证明体积。得益于 KZG 多项式承诺的高效压缩能力，Halo2 生成的证明极其紧凑。对于 2450 步的复杂证明，仅需 1.19 MB，而 ZKWASM 则需要 1344.38 MB（约 1.3 GB）。三个数量级的巨大缩减意味着我们的证明可以轻松嵌入区块链交易或在低带宽网络中传输，具有极高的工程实用价值。

\begin{table}[htbp]
\centering
\bicaption{证明规模对比：Halo2 vs ZKWASM}{UNSAT Case Proof Size Comparison}
\label{tab:proof-size}
\begin{tabular}{cccc}
\toprule
\multirow{2}{*}{\textbf{归结长度}} & \multirow{2}{*}{\textbf{子句宽度}} & \multicolumn{2}{c}{\textbf{证明大小}} \\
\cline{3-4}
 & & \textbf{ZKWASM} & \textbf{Halo2 (Ours)} \\
\midrule
20 & 5 & 9.21 MB & 32.50 KB \\
274 & 34 & 210.04 MB & 90.60 KB \\
1034 & 84 & 415.46 MB & 154.09 KB \\
1511 & 98 & 506.80 MB & 411.43 KB \\
2450 & 142 & 1344.38 MB & 1.19 MB \\
\bottomrule
\end{tabular}
\end{table}
